% SOURCE_AUTHORITY: sga5_fr_workpass.tex, lines 12273--12291 (LF-normalized unit).
% SOURCE_SHA256: 791F4EFFC5E02832D5D77ED03518C8156D6F07E4C8238B03545DB93D883FBB28
\subsection*{2. Donde se demuestra que cierto complejo de $\Lambda[G]$-módulos es perfecto}

Recordemos el resultado siguiente de SGA 4 XVII.

\begin{statement}{Lema 2.1}
Sean $X$ e $Y$ dos esquemas, y sea $f:X\to Y$ un morfismo propio, con $Y$ cuasicompacto. Denotemos por $A$ un anillo de torsión, y sea $F^\bullet$ un complejo de haces de $A$-módulos sobre $X$ de dimensión Tor finita (SGA 4 XVII); entonces $Rf_*(F^\bullet)$ tiene dimensión Tor finita.
\end{statement}

\begin{statement}{Proposición 2.2}
Sean $X$ un esquema propio sobre el cuerpo separablemente cerrado $k$, $G$ un grupo finito que actúa sobre $X$, $V$ un abierto de $X$ estable bajo $G$ y tal que $G$ actúa de manera admisible (SGA 1 V 1.7) y libremente sobre $V$, $\Lambda=\mathbb Z/\ell^n\mathbb Z$, con $\ell$ un número primo distinto de la característica de $k$, y $\Lambda_V$ el haz constante sobre $V$ asociado a $\Lambda$, dotado de su estructura trivial de haz con grupo de operadores $G$. Denotemos por
\[
i:V\longrightarrow X
\]
la inmersión de $V$ en $X$, y sea $\Lambda_{V,X}=i_!(\Lambda_V)$. En estas condiciones, el complejo de $\Lambda[G]$-módulos
\[
R\Gamma_X(\Lambda_{V,X})
\]
es un complejo perfecto (VIII 7).
\end{statement}
